\section{Discussion and extensions}

The rate characterization in \cref{obj:thm:degree-frontier,obj:thm:bounded-outcome-degree-frontier} separates two sources of design difficulty. The first is the local Bernoulli score energy \(A_d\), which is determined by the assignment probability \(p\), the exposed interaction order, and the number of coordinates in one neighborhood. The second is the additional out-degree overlap charge \(d\), which enters because one assignment coordinate can appear in the scores of as many as \(d\) units under \cref{obj:ass:bounded-degree}. Together they yield the scale
\[
B^2\min\left\{1,\frac{dA_d}{n}\right\},
\]
with the equivalent exposed-order form
\[
B^2\min\left\{1,\frac{d\binom d{k_\star(d,\beta,p)}}{n}\right\}
\]
up to constants depending on \((\beta,p)\). The constants obtained from the block-energy comparison are interpreted for a fixed \((\beta,p)\) with \(p\in(0,1)\). Under that fixed pair, the frontier is uniform in \(n\), \(d\), and \(B\).

For design planning, \(A_d\) is the local price of extracting the all-treated-versus-all-control contrast from Bernoulli variation inside one neighborhood, while \(dA_d/n\) is the population-level burden after neighborhood overlap. The linear branch \(dA_d/n\ll1\) is the regime where unprojected SNIPE variance is informative. In the stated bounded classes, when \(dA_d/n\) is bounded away from zero at the saturation scale, the minimax mean squared error remains of order \(B^2\) up to constants depending on \((\beta,p)\).

The energy \(A_d\) gives the local price of recovering the all-treated-versus-all-control contrast from Bernoulli variation within a neighborhood. Its definition in \cref{obj:def:block-score-energy} sums the squared Bernoulli contrast coefficients across raw monomials up to the available order \(\bar\beta_d\). The comparison in \cref{obj:thm:bounded-outcome-degree-frontier} shows that, for fixed \((\beta,p)\), this energy has the same order as \(\binom d{k_\star(d,\beta,p)}\). Thus the exposed nonzero contrast order determines how the local polynomial dimension enters the risk, while the assignment probability determines which orders are visible through the coefficients \(\Delta_r(p)\).

The overlap factor has a different interpretation. It is a graph-combinatorial charge rather than a local representer charge. The covariance terms in SNIPE are organized by common assignment coordinates across neighborhoods, and \cref{obj:thm:degree-frontier} bounds the number of such overlaps by \(d\binom d r\) at each order \(r\). This is the point at which bounded out-degree matters for mean squared error: it controls how many unit-level score contributions can share a given monomial. Related local-interference analyses also emphasize neighborhood growth and overlap as central design features \citep{Savje2021,Leung2022,Hu2022,Gao2025}; the characterization here converts that feature into the finite-population minimax factor \(dA_d/n\) for known bounded-degree low-order polynomial interference.

Complete directed blocks provide the sharp calibration for the same expression. In the block construction of \cref{obj:def:block-family}, each active unit has the full \(d\)-coordinate neighborhood, and the normalized representer direction produces least-favourable perturbations whose separation and affinity are governed by \(A_d/m\). When \(d\mid n\), \cref{obj:thm:bounded-outcome-degree-frontier} gives the exact unprojected SNIPE worst-case risk \(B^2A_d/m=B^2dA_d/n\) on complete blocks. The complete-block calculation matches the global upper bound's degree dependence and gives the exact unprojected-SNIPE worst-case risk on the block graph.

The local linear benchmark in \cref{obj:thm:sharp-local-linear-constant-and-representers} gives a complementary finite-dimensional interpretation. On fixed complete directed block graphs under the coefficient-mass schedule class, the minimax risk over block-local unbiased linear estimators is exactly \(B^2A_{d_t}/m_t\). This exact constant characterizes the block-local unbiased linear class on fixed complete directed blocks under the coefficient-mass schedule class: it optimizes over weights satisfying the block-local moment equations against every nonempty raw monomial up to \(\bar\beta_{d_t}\). The clipped estimator, which uses projection, is the procedure that attains the saturated branch of the bounded-class frontier. The criterion \(\Psi_{b,t}(w)\) identifies the blockwise extremal quantity that any proposed local linear weight must control, and the representer condition gives a sufficient route to the same benchmark through average \(L^2(P_Z)\) closeness to the complete-block SNIPE score. This places SNIPE within a broader local Riesz geometry while preserving the exact blockwise minimax constant.

The comparison between the coefficient-mass and uniformly bounded-outcome classes in \cref{obj:thm:bounded-outcome-degree-frontier} shows that the same degree dependence governs both boundedness formulations. The bounded-outcome class is larger, but the clipped SNIPE estimator projected to \([-2B,2B]\) attains the same \(B^2\min\{1,dA_d/n\}\) scale. This matters for applications in which the primitive restriction is naturally stated as a uniform potential-outcome envelope rather than an \(\ell_1\) coefficient-mass envelope: the design-based rate remains calibrated by the same local energy and overlap factor.

The fair-coin specialization in \cref{obj:thm:fair-coin-energy-frontier} makes the role of assignment probability especially transparent. At \(p=1/2\), the even-order Bernoulli contrasts vanish and the energy is the sum over odd orders. For first-order interference, \(A_d=4d\), so the minimax scale becomes \(B^2\min\{1,d^2/n\}\), with exact complete-block unprojected SNIPE risk \(4B^2d^2/n\) when \(d\mid n\). This gives a simple known-graph Bernoulli benchmark for comparisons with clustered assignment, unknown interference, and covariate-adjusted designs \citep{EichhornKhanUganderYu2026,WangLi2026}.

For fixed \((\beta,p)\) sequences with eventual positive degree and stable exposed order \(k_\star\), the comparison \(dA_d\asymp_{\beta,p} d\binom d{k_\star}\) identifies a sufficient and rate-equivalent vanishing condition for the displayed minimax scale. The heuristic condition is \(d\binom d{k_\star}=o(n)\). Under that fixed-pair interpretation, this corresponds to degree growth on the order \(d=o(n^{1/(k_\star+1)})\). For fair-coin assignment with first-order interactions \((\beta=1)\), the condition reads \(d=o(\sqrt n)\). At the saturation scale, for fixed \((\beta,p)\), the frontier identifies \(B^2\) as the constant-order minimax scale.

\subsection{Limitations and future work}

The results establish design-based minimax mean squared error rates for known bounded-degree low-order polynomial interference under common-probability Bernoulli assignment, together with exact complete-block constants for the stated local linear class.

The framework treats potential outcomes as exact low-order polynomials of the assignment vector, so the reported risk is assignment-design mean squared error for the fixed schedule. Open directions include additive idiosyncratic outcome noise and the associated variance floor, exact leading constants over all measurable estimators beyond the complete-block and local linear benchmarks, heterogeneous treatment probabilities, variance estimation and studentization for the SNIPE family, limit laws under growing-degree regimes, unknown or partially observed graphs, clustered assignment, and covariate adjustment. These extensions would connect the benchmark developed here to settings studied in recent work on local interference, graph uncertainty, and adjusted experimental estimators \citep{Savje2021,Leung2022,Hu2022,Gao2025,EichhornKhanUganderYu2026,WangLi2026}.
